% ===========================================================================
%  Jàng Afrig — Plan de Cours : Fondamentaux du Web (HTML, CSS, JavaScript)
%  Licence 1 DUT1-Informatique — ESP/UCAD — 2025–2026
%  Dr. El Hadji Bassirou TOURÉ
%  Compilé avec : xelatex (2–3 passes)
%  Version : courte (distribuable aux étudiants)
% ===========================================================================
\documentclass[11pt,a4paper]{article}

\usepackage{fontspec}
\setmainfont{Latin Modern Roman}
\setsansfont{Latin Modern Sans}
\setmonofont{Latin Modern Mono}
\usepackage{polyglossia}
\setdefaultlanguage{french}

\usepackage[margin=2.5cm]{geometry}
\usepackage{amsmath,amssymb}
\usepackage{graphicx}
\graphicspath{{figures/}}
\usepackage{xcolor}
\usepackage{booktabs}
\usepackage{hyperref}
\usepackage{tcolorbox}
\tcbuselibrary{theorems,breakable,skins}
\usepackage{enumitem}
\usepackage{fancyhdr}
\usepackage{array}
\usepackage{pifont}
\usepackage{longtable}
\usepackage{multirow}

\newcommand{\cmark}{\ding{51}}
\newcommand{\xmark}{\ding{55}}

\hypersetup{
    colorlinks=true,
    linkcolor=blue!60!black,
    urlcolor=blue!60!black,
    citecolor=blue!60!black
}

\newtcolorbox{objectifbox}[1][]{colback=blue!5,colframe=blue!50,fonttitle=\bfseries,title=#1,breakable}
\newtcolorbox{infobox}[1][]{colback=green!5,colframe=green!50!black,fonttitle=\bfseries,title=#1,breakable}
\newtcolorbox{attentionbox}[1][]{colback=orange!5,colframe=orange!70,fonttitle=\bfseries,title=#1,breakable}
\newtcolorbox{retenirbox}[1][]{colback=yellow!10,colframe=yellow!50!black,fonttitle=\bfseries,title=#1,breakable}
\newtcolorbox{prereqbox}[1][]{colback=cyan!5,colframe=cyan!50!black,fonttitle=\bfseries,title=#1,breakable}
\newtcolorbox{projetbox}[1][]{colback=teal!5,colframe=teal!50!black,fonttitle=\bfseries,title=#1,breakable}

\usepackage{listings}
\lstset{
  basicstyle=\ttfamily\small,
  backgroundcolor=\color{gray!8},
  frame=single,
  framerule=0.5pt,
  rulecolor=\color{gray!40},
  breaklines=true,
  breakatwhitespace=true,
  showstringspaces=false,
  tabsize=2,
  xleftmargin=4pt,
  xrightmargin=4pt,
  aboveskip=8pt,
  belowskip=8pt,
  literate={é}{{\'e}}1 {è}{{\`e}}1 {ê}{{\^e}}1 {à}{{\`a}}1 {ù}{{\`u}}1
           {ô}{{\^o}}1 {î}{{\^i}}1 {ç}{{\c{c}}}1 {É}{{\'E}}1
           {→}{{$\rightarrow$}}1
}

\pagestyle{fancy}
\fancyhf{}
\fancyhead[L]{\small Jàng Afrig --- Plan de Cours}
\fancyhead[R]{\small UCAD --- L1 DUT1-Info}
\fancyfoot[C]{\thepage}
\fancyfoot[L]{\small Dr. El Hadji Bassirou TOURÉ}
\fancyfoot[R]{\small 2025--2026}
\renewcommand{\headrulewidth}{0.4pt}
\renewcommand{\footrulewidth}{0.4pt}

\title{%
\vspace{-1cm}
\textbf{\LARGE Fondamentaux du Web}\\
\textbf{\LARGE HTML, CSS, JavaScript}\\[0.8cm]
\textbf{\huge Plan de Cours}\\[0.3cm]
\textbf{\Large Projet fil rouge : Jàng Afrig}\\[0.5cm]
\large Licence 1 DUT1-Informatique --- Semestre 2\\
2025--2026}
\author{Dr. El Hadji Bassirou TOURÉ\\
Département Génie Informatique\\
Université Cheikh Anta Diop de Dakar}
\date{}

\begin{document}

\maketitle
\thispagestyle{empty}

\vspace{0.5cm}

% ===========================================================================
% BOÎTE PRÉREQUIS
% ===========================================================================
\begin{prereqbox}[Informations pratiques]
\begin{center}
\begin{tabular}{ll}
\toprule
\textbf{Volume horaire} & 15 séances de 2h = \textbf{30h} \\
\textbf{Public} & Licence 1 DUT1-Informatique --- DGI / ESP / UCAD \\
\textbf{Semestre} & Semestre 2 (2\textsuperscript{e} semestre de L1) \\
\textbf{Prérequis} & Algorithmique du S1. Savoir utiliser un ordinateur. \\
\textbf{Outils} & VS Code, Live Server, Git, GitHub, Chrome ou Firefox \\
\textbf{Projet} & \textbf{Jàng Afrig} --- quiz géographique sur les 54 pays africains \\
\textbf{Évaluation} & Participation + avancement Git (voir détail p.~\pageref{sec:evaluation}) \\
\bottomrule
\end{tabular}
\end{center}
\end{prereqbox}

\begin{objectifbox}[Ce que vous saurez faire à la fin du cours]
\begin{enumerate}[nosep]
    \item Écrire une page web complète en HTML5 (structure, sémantique, formulaires, tableaux)
    \item Mettre en forme une page avec CSS3 (sélecteurs, box model, Flexbox, responsive mobile)
    \item Programmer en JavaScript moderne (variables, fonctions, tableaux, objets, \texttt{forEach}, \texttt{map}, \texttt{filter})
    \item Modifier le contenu d'une page depuis JavaScript (DOM et événements)
    \item Récupérer des données d'une API distante avec \texttt{fetch} et \texttt{async/await}
    \item Sauvegarder des données dans le navigateur (\texttt{localStorage})
    \item Versionner votre code avec Git et le partager sur GitHub
    \item Déployer une application web en ligne sur GitHub Pages
\end{enumerate}
\end{objectifbox}

\tableofcontents

\newpage

% ===========================================================================
\section{Présentation du cours}
% ===========================================================================

Ce cours est votre \textbf{premier contact avec un vrai langage de
programmation qui s'exécute dans un outil réel}. Vous avez vu l'algorithmique
en pseudo-code au S1 ; vous allez maintenant écrire du code qui
\textit{marche}, dans un navigateur, et qui produit des résultats visibles.

Au terme du semestre, vous aurez construit \textbf{Jàng Afrig}, un quiz
géographique sur les 54 pays africains, accessible depuis n'importe quel
téléphone via une URL publique que vous pourrez partager à votre famille.

\begin{infobox}[Pourquoi 100\% vanilla ?]
Vous entendrez parler de frameworks (React, Vue, Next.js\ldots). Ce sont des
outils très puissants, mais ils \textbf{cachent} ce qui se passe vraiment
dans le navigateur. Avant d'utiliser un framework, il faut comprendre ce
qu'il fait à votre place. Ce cours se concentre donc sur les briques
standard du Web --- pas de framework, pas de \texttt{npm}. Les frameworks
viendront en L2 (cours PRC) et L3 (cours IPDL1).
\end{infobox}

\subsection{Deux types de séances : CM et Lab}

Le cours alterne deux types de séances, qui se complètent :

\begin{itemize}[leftmargin=*]
    \item \textbf{CM (Cours Magistral).} 4 séances au total. On introduit
          les concepts avec des analogies, des exemples visuels, une démo
          en direct, et un mini-exercice en classe. Aucun devoir à la
          maison après un CM.

    \item \textbf{Lab (TP guidé).} 10 séances au total. On met les mains
          dans le code, on construit une partie de Jàng Afrig, on fait des
          exercices d'application, et on termine par la routine Git.
\end{itemize}

La dernière séance (S15) est un \textbf{hybride} : première heure sur le
déploiement en ligne, deuxième heure démo collective.

\subsection{Déroulement d'une séance}

\begin{center}
\begin{tabular}{@{}p{2.5cm}p{5.5cm}p{5.5cm}@{}}
\toprule
\textbf{Moment} & \textbf{CM (2h)} & \textbf{Lab (2h)} \\
\midrule
5 min    & Rappel séance précédente & Rappel séance précédente \\
10 min   & Analogie + question centrale & Démo de l'objectif du jour \\
60 min   & Contenu théorique en 3-4 blocs & Vous codez en suivant le pas-à-pas \\
30 min   & Mini-exercice en classe & Exercices d'application + avancement Jàng Afrig \\
15 min   & Résumé + teaser du Lab suivant & Routine Git + annonce S suivante \\
\bottomrule
\end{tabular}
\end{center}

\newpage

% ===========================================================================
\section{Le projet fil rouge : Jàng Afrig}
% ===========================================================================

\textbf{Jàng} signifie \textit{«\,apprendre, étudier\,»} en wolof, et \textbf{Afrig} désigne l'Afrique. Le nom évoque donc littéralement \textit{«\,apprendre l'Afrique\,»} --- ce que fait l'étudiant en jouant au quiz tout en construisant l'application.
Jàng Afrig est une application web de \textbf{quiz géographique} sur les 54
pays africains. L'utilisateur choisit un mode de jeu, répond à 10 questions
à choix multiples, obtient un score, et garde un classement local de ses
meilleures performances.

\subsection{Les modes de jeu}

\begin{center}
\begin{tabular}{ll}
\toprule
\textbf{Mode} & \textbf{Question type} \\
\midrule
Capitales      & «\,Quelle est la capitale du Mali ?\,» \\
Drapeaux       & «\,À quel pays appartient ce drapeau ?\,» \\
Langues        & «\,Dans quel pays parle-t-on le swahili ?\,» \\
Voisins        & «\,Combien de voisins a l'Algérie ?\,» \\
Chrono         & Maximum de bonnes réponses en 60 secondes \\
\bottomrule
\end{tabular}
\end{center}

\subsection{Comment on le construit}

\textbf{Jàng Afrig n'est pas un projet à rendre.} C'est l'application qu'on
construit \textbf{ensemble}, séance après séance, en Lab. À chaque Lab,
l'application gagne une fonctionnalité visible. Au bout de 15 séances,
c'est une vraie application en ligne.

\begin{projetbox}[Les données]
L'application utilisera l'API publique \href{https://restcountries.com}{REST
Countries v3.1}, qui fournit gratuitement des informations sur tous les pays
du monde : nom, capitale, drapeau, langues, frontières, superficie. Aucune
inscription, aucune clé API, aucune limite.

Un seul appel au démarrage, puis toutes les questions sont générées
localement à partir du tableau des 54 pays africains.
\end{projetbox}

\subsection{Les pages de l'application}

\begin{enumerate}[nosep]
    \item \textbf{Accueil} : choix du mode + saisie du pseudo + meilleur score
    \item \textbf{Quiz} : énoncé + 4 propositions + barre de progression
    \item \textbf{Résultat} : score final + bouton rejouer
    \item \textbf{Mes Scores} : classement local par mode (top 5)
    \item \textbf{À propos} : sources des données, technologies
\end{enumerate}

\newpage

% ===========================================================================
\section{Outils à installer}
% ===========================================================================

\textbf{Avant la première séance}, vous devez avoir installé tous les
outils et fait votre premier cycle Git. Une fiche dédiée \texttt{Lab0\_Setup.pdf}
vous guide pas à pas, avec captures d'écran. \textbf{Faites-la à la maison
la veille du cours}.

\begin{center}
\begin{tabular}{lll}
\toprule
\textbf{Outil} & \textbf{Usage} & \textbf{Lien} \\
\midrule
\textbf{VS Code} & Éditeur de code & \href{https://code.visualstudio.com}{code.visualstudio.com} \\
\textbf{Live Server} & Extension VS Code (1 clic) & Marketplace VS Code \\
\textbf{Git} & Suivre les modifications du code & \href{https://git-scm.com}{git-scm.com} \\
\textbf{Compte GitHub} & Hébergement du code en ligne & \href{https://github.com}{github.com} \\
\textbf{Chrome ou Firefox} & Navigateur + outils de développement & Préinstallé \\
\bottomrule
\end{tabular}
\end{center}

\begin{attentionbox}[Aucun npm, aucun framework]
Le cours est volontairement vanilla. \textbf{Aucune installation Node.js},
aucun \texttt{package.json}, aucun \texttt{npm install}.
\end{attentionbox}

\subsection{Git : la routine de fin de Lab}

À la fin de chaque Lab, vous exécuterez la même séquence rituelle :

\begin{lstlisting}
git add .
git commit -m "v0.X : titre du Lab"
git push
git tag v0.X
git push --tags
\end{lstlisting}

Cette routine devient automatique au bout de 3-4 Labs. Les commandes Git
elles-mêmes sont toutes enseignées dans le Lab 0 --- \textbf{aucun prérequis}
n'est demandé.

\newpage

% ===========================================================================
\section{Plan des 15 séances}
% ===========================================================================

\begin{center}
\begin{tabular}{clllc}
\toprule
\textbf{N\textsuperscript{o}} & \textbf{Type} & \textbf{Bloc} & \textbf{Titre} & \textbf{Tag Git} \\
\midrule
\textbf{Lab 0} & Lab & Setup & Mise en route (à la maison) & \texttt{v0} \\
\midrule
S1   & \textbf{CM1}   & HTML & Le Web et HTML                                      & --- \\
S2   & Lab 1          & HTML & Ma première page HTML                               & \texttt{v0.1} \\
S3   & Lab 2          & HTML & Sémantique, formulaires, tableaux                   & \texttt{v0.2} \\
\midrule
S4   & \textbf{CM2}   & CSS  & Sélecteurs, cascade, box model, Flexbox             & --- \\
S5   & Lab 3          & CSS  & CSS en pratique : styliser Jàng Afrig                & \texttt{v0.3} \\
S6   & Lab 4          & CSS  & Layouts + responsive mobile                          & \texttt{v0.4} \\
\midrule
S7   & \textbf{CM3}   & JS   & JavaScript : le langage qui anime les pages         & --- \\
S8   & Lab 5          & JS   & JS pratique : variables, fonctions, conditions      & \texttt{v0.5} \\
S9   & Lab 6          & JS   & Tableaux et objets (\texttt{forEach}, \texttt{map}, \texttt{filter}) & \texttt{v0.6} \\
S10  & Lab 7          & JS   & DOM et événements : premier quiz fonctionnel        & \texttt{v0.7} \\
\midrule
S11  & \textbf{CM4}   & API  & Le Web moderne : APIs, JSON, asynchrone             & --- \\
S12  & Lab 8          & API  & Fetch + \texttt{async/await} : brancher REST Countries & \texttt{v0.8} \\
S13  & Lab 9          & API  & Tous les modes de jeu avec l'API                    & \texttt{v0.9} \\
S14  & Lab 10         & JS   & \texttt{localStorage}, chrono, polish               & \texttt{v0.14} \\
\midrule
\textbf{S15} & \textbf{Hybride} & Final & \textbf{Déploiement GitHub Pages + démo collective} & \textbf{\texttt{v1.0}} \\
\bottomrule
\end{tabular}
\end{center}

\bigskip

\begin{infobox}[Le fil narratif]
Chaque bloc est construit autour d'un CM qui pose les concepts, suivi de
Labs qui les appliquent sur Jàng Afrig. À chaque fin de Lab, un tag Git
\texttt{v0.X} marque la progression. À la fin, le tag \texttt{v1.0} marque
la version publiée en ligne.
\end{infobox}

\bigskip

\begin{retenirbox}[À chaque CM, chaque Lab, un PDF]
Pour chaque séance, un document PDF dédié vous est distribué :
\begin{itemize}[nosep]
    \item \textbf{CMX\_[titre].pdf} pour les cours magistraux
    \item \textbf{LabX\_[titre].pdf} pour les TP guidés
\end{itemize}
Le prof projette le PDF en direct, vous le suivez sur votre écran ou en
papier. Tout est écrit, rien n'est à noter à la volée.
\end{retenirbox}

\newpage

% ===========================================================================
\section{Évaluation}
\label{sec:evaluation}
% ===========================================================================

\subsection{Philosophie}

Ce cours \textbf{ne comporte pas de projet noté final} à rendre. Le projet
Jàng Afrig n'est pas un livrable séparé : c'est l'application qu'on construit
ensemble en Labs, séance après séance. L'évaluation porte sur la
\textbf{participation continue} et la \textbf{progression visible}.

\subsection{Critères d'évaluation}

\begin{center}
\begin{tabular}{lp{8cm}c}
\toprule
\textbf{Critère} & \textbf{Comment c'est observé} & \textbf{Poids} \\
\midrule
Présence et participation active & Présence aux séances + interactions en Labs & 20\% \\
Avancement Git & Tags posés de \texttt{v0.1} à \texttt{v1.0} & 30\% \\
Qualité du code & Code lisible, indenté, commenté & 20\% \\
Fonctionnalité de l'app & L'app finale fonctionne sur le téléphone du jury & 20\% \\
Déploiement réussi & URL GitHub Pages accessible en S15 & 10\% \\
\bottomrule
\end{tabular}
\end{center}

\begin{retenirbox}[Le seul vrai test]
À la fin du cours, vous devez être capable de :
\begin{enumerate}[leftmargin=*,nosep]
  \item Ouvrir un terminal et faire un commit Git
  \item Écrire une page HTML sémantique
  \item Styliser une page avec Flexbox responsive
  \item Manipuler le DOM en JavaScript
  \item Manipuler un tableau avec \texttt{forEach}, \texttt{map}, \texttt{filter}
  \item Récupérer des données d'une API avec \texttt{async}/\texttt{await}
  \item Persister des données avec \texttt{localStorage}
  \item Déployer un site statique sur GitHub Pages
\end{enumerate}
Si oui, le cours a réussi.
\end{retenirbox}

\newpage

% ===========================================================================
\section{Ressources}
% ===========================================================================

\subsection{Pendant le cours}

\begin{itemize}[leftmargin=*,nosep]
    \item \textbf{MDN Web Docs} --- la documentation de référence.\\
          \url{https://developer.mozilla.org/fr}
    \item \textbf{REST Countries v3.1} --- l'API du projet.\\
          \url{https://restcountries.com}
\end{itemize}

\subsection{Pour aller plus loin (optionnel)}

\begin{itemize}[leftmargin=*,nosep]
    \item Marijn Haverbeke, \textit{Eloquent JavaScript} (3\textsuperscript{e} éd.), gratuit en ligne : \url{https://eloquentjavascript.net}
    \item Stanford CS193X --- \textit{Web Programming Fundamentals} (cours d'inspiration) : \url{https://web.stanford.edu/class/archive/cs/cs193x/cs193x.1176}
\end{itemize}

\subsection{Cours connexes du cursus}

\begin{itemize}[leftmargin=*,nosep]
    \item \textbf{L2 S4 --- Programmation par Réutilisation de Composants (PRC).}
          Suite naturelle : React, Flask, fil rouge SénTransport.
    \item \textbf{L2 S4 --- Intégration de Modèles IA (IMA).} APIs, Docker, fil rouge SénSanté.
    \item \textbf{L3 --- IPDL1.} Ingénierie des processus, Next.js, équipes de 6.
\end{itemize}

\vspace{1.5cm}

\begin{center}
\fbox{
\begin{minipage}{0.85\textwidth}
\centering
\vspace{0.5cm}
\textbf{\Large Bon semestre à toutes et tous \cmark}\\[0.3cm]
\normalsize
\textit{Le code, ça s'apprend en codant.}\\[0.3cm]
\vspace{0.3cm}
\end{minipage}
}
\end{center}

\end{document}
